\subsection{The definiteness on \texorpdfstring{$T_{2g}$}{T2g}}

With the standard cube orientation fixed above, let
\[
A:=\{1,2,7,8\},\qquad B:=\{3,4,5,6\},
\qquad
G_A:=\sum_{i\in A}G_{p_i},\qquad G_B:=\sum_{i\in B}G_{p_i}.
\]
The \(T_{2g}\)-generator from the decomposition section is
\[
v_{T_{2g}}
=
\sum_{i\in A}\frac{\partial_\phi|_{p_i}}{\sin\theta_0}
-
\sum_{i\in B}\frac{\partial_\phi|_{p_i}}{\sin\theta_0}.
\]
Since \(\Hess\) is invariant under the cube symmetries and \(T_{2g}\) is irreducible, it is enough to determine the sign of \(\Hess(v_{T_{2g}},v_{T_{2g}})\). By the general Hessian formula obtained earlier, this sign is opposite to the sign of
\[
I_{T_{2g}}
:=
\int_{\uS^2}\partial_\phi G_A\,\partial_\phi G_B\,\dd(\mu^++\mu^-).
\]
Thus it remains to prove that \(I_{T_{2g}}<0\).

Let \(R_{-\pi/2}\) be the rotation \(\phi\mapsto\phi-\pi/2\) about the \(Z\)-axis. This is a symmetry of the cube, and for the two four-point configurations defining \(G_A\) and \(G_B\) one checks that
\[
G_B(q)=G_A(R_{-\pi/2}(q)).
\]
Because \(R_{-\pi/2}\) only shifts the azimuthal angle, it commutes with \(\partial_\phi\).

Now parametrize \(\uS^2\) by
\[
q(x,\psi):=\bigl(r\cos\psi,r\sin\psi,x\bigr),
\qquad
r:=\sqrt{1-x^2},
\qquad
-1<x<1,\quad 0\leq \psi<2\pi,
\]
so that \(x=\cos\theta\) and \(\psi=\phi\). Set
\[
h(x,\psi):=(\partial_\phi G_A)(q(x,\psi)).
\]
With
\[
a:=\sqrt{\frac23},
\qquad
b:=\frac1{\sqrt3},
\]
the explicit formula from \texttt{sol\_v2\_1\_formulas.tex} is
\[
h(x,\psi)
=
-\frac{x\,r\,\sin\psi\bigl(3+2r^2\cos^2\psi-x^2\bigr)}
{9\pi\,P_A(x,r\cos\psi)},
\]
where
\[
P_A(x,z):=
\bigl((z-a)^2+b^2(1-x^2-z^2)\bigr)
\bigl((z+a)^2+b^2(1-x^2-z^2)\bigr).
\]
Since \(R_{-\pi/2}(q(x,\psi))=q(x,\psi-\pi/2)\), we obtain
\[
\partial_\phi G_B(q(x,\psi))
=
\frac{\partial}{\partial\psi}\Bigl(G_A\bigl(q(x,\psi-\pi/2)\bigr)\Bigr)
=
h(x,\psi-\pi/2).
\]

Let \(w\) be the density of \(\mu^++\mu^-\) in these coordinates:
\[
\dd(\mu^++\mu^-)=w(x,\psi)\,\dd\psi\,\dd x.
\]
Therefore
\[
I_{T_{2g}}
=
\int_{-1}^{1}\int_{0}^{2\pi}
h(x,\psi)\,h(x,\psi-\pi/2)\,w(x,\psi)\,\dd\psi\,\dd x
=:K_1.
\]
This is exactly the quantity \(K_1\) from the lemma of the previous subsection, with \(f=h\) and \(\rho=w\). We now verify its hypotheses.

\begin{enumerate}
\item The function \(h\) is odd in \(\psi\):
\[
h(x,-\psi)=-h(x,\psi).
\]
Indeed, \(\sin(-\psi)=-\sin\psi\), while \(P_A(x,z)\) is even in \(z\), hence \(P_A(x,r\cos(-\psi))=P_A(x,r\cos\psi)\).

\item The strict sign condition
\[
h(x,\psi)<0
\qquad
\text{for all }x\in(0,1)\text{ and }\psi\in(0,\pi)
\]
is proved in \texttt{sol\_v2\_5\_sign.tex}.

\item For each fixed \(x\in(0,1)\), the map \(\psi\mapsto h(x,\psi)\) is strictly increasing on \((0,\pi)\). This is proved in \texttt{sol\_v2\_3\_monotonicity.tex}.

\item Since \(\mu^++\mu^-\) is a positive measure, \(w\geq 0\). Moreover, the \(C_4\)-symmetry of the cube about the \(Z\)-axis gives
\[
w(x,\psi+\pi/2)=w(x,\psi),
\]
and the equatorial symmetry \((X,Y,Z)\mapsto(X,Y,-Z)\) gives
\[
w(-x,\psi)=w(x,\psi).
\]
\end{enumerate}

Hence the lemma yields \(K_1<0\). Consequently,
\[
\int_{\uS^2}\partial_\phi G_A\,\partial_\phi G_B\,\dd(\mu^++\mu^-)<0.
\]
Therefore \(\Hess(v_{T_{2g}},v_{T_{2g}})>0\), and since \(\Hess|_{T_{2g}}\) is a scalar on the irreducible \(T_{2g}\)-block, the Hessian is positive definite on \(T_{2g}\).